\begin{textsample}{NEG Dim 1 – 1950 – Score -2.00 – 1950/tv\_com\_1950\_40.txt}  \label{ex:f1_neg_014}
The commercial opens in black and white on a crowded kitchen counter and sink area.
Stacks of dirty dishes fill the space : a pile of plates sits toward the back, a large serving platter lies to the right, cups and saucers are clustered on the side, and several drinking glasses stand in front.
Many pieces look smeared or cloudy, with food residue or stains visible on the plates and glassware.
Manicured hands enter from the right and add more cups and saucers to the already crowded counter, emphasizing the size of the dishwashing job.
A hand reaches toward the sink faucet.
Beside the faucet, a dark cylindrical can of Lux Liquid Detergent is set upright on the counter.
The label is bold and high-contrast, with very large white letters spelling “ LUX ” and smaller text below reading “ LIQUID DETERGENT. ” The can has a metal top with a short raised opening or cap-like spout.
The hand steadies the can near the sink, turning it slightly so the label faces forward.
The can is held close to the sink and tilted over a spoon.
A thick, dark liquid detergent pours from the small opening into the spoon, filling it.
The spoon is held carefully to the right of the can, and the product flows in a controlled stream.
After the spoonful is shown, the spoon is tipped and the liquid runs off toward the sink.
Water begins running from the curved faucet into the basin, and \textbf{suds} quickly build up, spreading into a frothy layer across the bottom of the sink.
The detergent can briefly appears amid the growing foam, then stands prominently again near the faucet with its top opening visible.
Attention moves to the can ’s top.
A manicured finger demonstrates the small metal pouring opening and the raised rim around it, touching and wiping around the spout area.
The hand moves across the top and side of the can, showing the opening and the surrounding metal surface cleanly, as if emphasizing that the container can be handled without mess.
The can remains upright beside the faucet, the large “ LUX ” name facing outward.
The sink is now full of \textbf{suds}.
The Lux can stands to the left of the faucet while clean water continues to run.
A hand reaches into the foam and lifts a drinking glass.
In close view, the glass is washed under the running water and in the \textbf{suds}.
The hand rotates it, rubbing around the outside and inside.
The glass, initially cloudy and wet, becomes clearer as water flows over it.
The fingers are slender with dark polished nails, and the glass is turned repeatedly so the rim, sides, and base are rinsed.
The demonstration shifts to a large plate or platter.
A hand holds it in the sudsy sink under the running faucet.
Dark streaks and greasy-looking marks are visible around the rim and surface.
The hand rubs and swirls the plate through the foam, turning it so water and \textbf{suds} wash across the surface.
As the plate is moved back and forth, the dark residue breaks up and disappears from the central area, leaving a broad clean circle.
The remaining marks around the rim are worked over with the fingers and water until the plate appears clean and glossy.
The plate is lifted and angled toward the viewer, its smooth pale center framed by a decorative rim.
The wider sink area returns.
The Lux can remains at the left near the faucet, with stacks of dishes and glasses in the background.
The water keeps running into the \textbf{foamy} basin.
The freshly washed plate is lowered out of view.
A hand hovers over the \textbf{suds}, then the two hands are shown together above the sink, rubbing palm to palm and across the fingers as if to show how the detergent feels on the skin.
The hands are wet and sudsy, but they appear relaxed and unhurried.
The product can appears again in close view.
The label “ LUX ” and “ LIQUID DETERGENT ” are clear on the front.
A hand points to the metal opening at the top and then moves away, leaving the raised spout visible.
The hand returns with the cap-like piece and places it back onto the top opening, pressing it down to close the container.
The can is steady in another hand, centered against a simple kitchen surface.
The final scene changes to a studio-like interior.
A middle-aged, light-skinned man with light wavy hair sits facing forward, wearing a light-colored suit jacket, dress shirt, and tie.
A small upright microphone stands in front of him at chest level.
Behind him are a lamp on the left and a softly patterned wall or curtain.
He holds the Lux Liquid Detergent can beside his face so the large “ LUX ” label is visible.
As he speaks, he shifts the can slightly, looks toward the viewer, and gestures with his free hand.
He points toward the can and taps or indicates the label while keeping it upright.
Near the end, he raises the can a little, presenting it directly, with the capped spout visible on top.

% matched lemmas: foamy, suds
\end{textsample}
